\documentclass{article}

% packages
\usepackage{amsmath, amsthm, thmtools, amsfonts, amssymb, luacode, catchfile, tikzducks, hyperref, ifthen}
\ifcsname c@kobocompile\endcsname
	\usepackage[a5paper, total={1072pt, 1448pt}, margin=10pt, includeheadfoot]{geometry} % set page margins
\else
	\usepackage[a4paper, margin=50pt, includeheadfoot]{geometry}
\fi
\usepackage[shortlabels]{enumitem}
\usepackage[skip=3pt, indent=0pt]{parskip}

% language
\usepackage[bidi=basic, layout=tabular, provide=*]{babel}
\ifcsname c@english\endcsname
	\babelprovide[main, import]{english}
\else
	\babelprovide[main, import]{hebrew}
	\babelprovide{rl}
\fi
%\babelfont{rm}{Libertinus Serif}
\babelfont{rm}[Renderer=Harfbuzz]{Libertinus Serif}
\babelfont{sf}{Libertinus Sans}
\babelfont{tt}{Libertinus Mono}

% style
\AddToHook{cmd/section/before}{\clearpage}	% Add line break before section
\linespread{1.3}
\setcounter{secnumdepth}{0}		% Remove default number tags from sections, this won't do well with theorems
\AtBeginDocument{\setlength{\belowdisplayskip}{3pt}}
\AtBeginDocument{\setlength{\abovedisplayskip}{3pt}}
\graphicspath{ {../images/} }

% operators
\DeclareMathOperator\cis{cis}
\DeclareMathOperator\Sp{Sp}
\DeclareMathOperator\tr{tr}
\DeclareMathOperator\im{Im}
\DeclareMathOperator\re{Re}
\DeclareMathOperator\diag{diag}
\DeclareMathOperator*\lowlim{\underline{lim}}
\DeclareMathOperator*\uplim{\overline{lim}}
\DeclareMathOperator\rng{rng}
\DeclareMathOperator\Sym{Sym}
\DeclareMathOperator\Arg{Arg}
\DeclareMathOperator\Log{Log}
\DeclareMathOperator\dom{dom}
\DeclareMathOperator\supp{Supp}
\DeclareMathOperator\var{Var}
\DeclareMathOperator\cov{Cov}

% commands
%\renewcommand\qedsymbol{\textbf{מש''ל}}
%\renewcommand\qedsymbol{\fbox{\emoji{lizard}}}
\newcommand{\Aa}[0]{\mathcal{A}}
\newcommand{\Bb}[0]{\mathcal{B}}
\newcommand{\CC}[0]{\mathbb{C}}
\newcommand{\Cc}[0]{\mathcal{C}}
\newcommand{\EE}[0]{\mathbb{E}}
\newcommand{\FF}[0]{\mathbb{F}}
\newcommand{\Ff}[0]{\mathcal{F}}
\newcommand{\Ii}[0]{\mathcal{I}}
\newcommand{\Gg}[0]{\mathcal{G}}
\newcommand{\Ll}[0]{\mathcal{L}}
\newcommand{\Mm}[0]{\mathcal{M}}
\newcommand{\NN}[0]{\mathbb{N}}
\newcommand{\Nn}[0]{\mathcal{N}}
\newcommand{\PP}[0]{\mathbb{P}}
\newcommand{\Pp}[0]{\mathcal{P}}
\newcommand{\QQ}[0]{\mathbb{Q}}
\newcommand{\RR}[0]{\mathbb{R}}
\newcommand{\Rr}[0]{\mathcal{R}}
\newcommand{\Ss}[0]{\mathcal{S}}
\newcommand{\TT}[0]{\mathbb{T}}
\newcommand{\Uu}[0]{\mathcal{U}}
\newcommand{\Vv}[0]{\mathcal{V}}
\newcommand{\Ww}[0]{\mathcal{W}}
\newcommand{\ZZ}[0]{\mathbb{Z}}
\newcommand{\acts}[0]{\circlearrowright}
\newcommand{\explain}[2] {
	\begin{flalign*}
		 && \text{#2} && \text{#1}
	\end{flalign*}
}
\newcommand{\maketitleprint}[0]{ \begin{center}
	%\begin{tikzpicture}[scale=3]
	%	\duck[graduate=gray!20!black, tassel=red!70!black]
	%\end{tikzpicture}	
	\includegraphics[width=6cm]{cover}
\end{center}
}

% theorem commands
\newtheoremstyle{c_remark}
	{}	% Space above
	{}	% Space below
	{}% Body font
	{}	% Indent amount
	{\bfseries}	% Theorem head font
	{}	% Punctuation after theorem head
	{.5em}	% Space after theorem head
	{\thmname{#1}\thmnumber{ #2}\thmnote{ \normalfont{\text{(#3)}}}}	% head content
\newtheoremstyle{c_definition}
	{3pt}	% Space above
	{3pt}	% Space below
	{}% Body font
	{}	% Indent amount
	{\bfseries}	% Theorem head font
	{}	% Punctuation after theorem head
	{.5em}	% Space after theorem head
	{\thmname{#1}\thmnumber{ #2}\thmnote{ \normalfont{\text{(#3)}}}}	% head content
\newtheoremstyle{c_plain}
	{3pt}	% Space above
	{3pt}	% Space below
	{\itshape}% Body font
	{}	% Indent amount
	{\bfseries}	% Theorem head font
	{}	% Punctuation after theorem head
	{.5em}	% Space after theorem head
	{\thmname{#1}\thmnumber{ #2}\thmnote{ \text{(#3)}}}	% head content

\ifcsname c@english\endcsname
	\theoremstyle{plain}
	\newtheorem{theorem}{Theorem}[section]
	\newtheorem{lemma}[theorem]{Lemma}
	\newtheorem{proposition}[theorem]{Proposition}
	\newtheorem*{proposition*}{Proposition}
	%\newtheorem{corollary}[theorem]{אין חלופה עברית}

	\theoremstyle{definition}
	\newtheorem{definition}[theorem]{Definition}
	\newtheorem*{definition*}{Definition}
	\newtheorem{example}{Example}[section]
	\newtheorem{exercise}{Exercise}[section]

	\theoremstyle{remark}
	\newtheorem*{remark}{Remark}
	\newtheorem*{solution}{Solution}
	\newtheorem{conclusion}[theorem]{Conclusion}
	\newtheorem{notation}[theorem]{Notation}
\else
	\theoremstyle{c_plain}
	\newtheorem{theorem}{משפט}[section]
	\newtheorem{lemma}[theorem]{למה}
	\newtheorem{proposition}[theorem]{טענה}
	\newtheorem*{proposition*}{טענה}
	%\newtheorem{corollary}[theorem]{אין חלופה עברית}

	\theoremstyle{c_definition}
	\newtheorem{definition}[theorem]{הגדרה}
	\newtheorem*{definition*}{הגדרה}
	\newtheorem{example}{דוגמה}[section]
	\newtheorem{exercise}{תרגיל}[section]

	\theoremstyle{c_remark}
	\newtheorem*{remark}{הערה}
	\newtheorem*{solution}{פתרון}
	\newtheorem{conclusion}[theorem]{מסקנה}
	\newtheorem{notation}[theorem]{סימון}
\fi

% Questions related commands
\newcounter{question}
\setcounter{question}{1}
\newcounter{sub_question}
\setcounter{sub_question}{1}

\ifcsname c@english\endcsname
	\newcommand{\question}[1][0]{
		\ifthenelse{#1 = 0}{}{\setcounter{question}{#1}}
		\section{Question \arabic{question}}
		\addtocounter{question}{1}
		\setcounter{sub_question}{1}
	}

	\newcommand{\subquestion}[1][0]{
		\ifthenelse{#1 = 0}{}{\setcounter{sub_question}{#1}}
		\subsection{Part \alph{sub_question}}
		\addtocounter{sub_question}{1}
	}
\else
	\newcommand{\question}[1][0]{
		\ifthenelse{#1 = 0}{}{\setcounter{question}{#1}}
		\section{שאלה \arabic{question}}
		\addtocounter{question}{1}
		\setcounter{sub_question}{1}
	}

	\newcommand{\subquestion}[1][0]{
		\ifthenelse{#1 = 0}{}{\setcounter{sub_question}{#1}}
		\subsection{סעיף \localecounter{letters.gershayim}{sub_question}}
		\addtocounter{sub_question}{1}
	}
\fi

% import lua and start of document
\directlua{common = require ('../common')}

\GetEnv{AUTHOR}

% headers
\author{\AUTHOR}
\date\today

\title{פתרון מטלה 06 --- פונקציות מרוכבות, 80519}

\begin{document}
\maketitle
\maketitleprint{}

\question{}
נחשב את האינטגרלים המסילתיים הנתונים.

\subquestion{}
\[
	\int_\gamma (2z - 3\overline{z} + 1)\ dz
\]
עבור $\gamma = 3 \cos t + 2i \sin t$ ב־$t \in [0, 2\pi]$.
\begin{solution}
	נתחיל בחישוב הכרחי:
	\[
		\gamma'(t) = -3 \sin t + 2i \cos t
	\]
	ונעבור לחישוב האינטגרל
	\begin{align*}
		\int_\gamma (2z - 3\overline{z} + 1)\ dz
		& = \int_0^{2\pi} (2(3 \cos t + 2i \sin t) - 3(3 \cos t - 2i \sin t) + 1)(-3 \sin t + 2i \cos t)\ dt \\
		& = \int_0^{2\pi} (-3 \cos t + 10i \sin t + 1)(-3 \sin t + 2i \cos t)\ dt \\
		& = \int_0^{2\pi} \frac{9}{2} \sin(2t) - 6 - 24i \sin^2 t - 10 \sin(2t) - 3 \sin t + 2i \cos t\ dt \\
		& = \int_0^{2\pi} - 6 - 6 {(e^{it} - e^{-it})}^2\ dt \\
		& = -6 \int_0^{2\pi} -1 + e^{2it} + e^{-2it}\ dt \\
		& = -6 {\left[ -t - \frac{i}{2} e^{2it} + \frac{i}{2} e^{-2it}\right]}_0^{2\pi}
	\end{align*}
	כאשר במעבר האחרון השתמשנו באינטגרל של פונקציה טריגונומטרית אפס בתחום.
\end{solution}

\subquestion{}
\[
	\int_\gamma \cos(\re(z))\ dz
\]
עבור $\gamma(t) = i + e^{it}$ עבור $t \in [-\pi, \pi]$.
\begin{solution}
	נובע $\gamma'(t) = i e^{it}$ ונבחין כי $\cos(\cos(t)) = -\cos(\cos(t))$ והפונקציה הזו זוגית, לכן
	\begin{align*}
		\int_\gamma \cos(\re(z))\ dz
		& = \int_{-\pi}^{\pi} \cos(\re(i + e^{it})) \cdot i e^{it}\ dt \\
		& = \int_{-\pi}^{\pi} \cos(\cos(t)) \cdot i e^{it}\ dt \\
		& = \int_{-\pi}^{\pi} \cos(\cos(t)) (i \cos t - \sin t)\ dt \\
		& = i \int_{-\pi}^{\pi} \cos(\cos(t)) \cos t\ dt - \int_{-\pi}^{\pi} \cos(\cos(t)) \sin t\ dt \\
		& = i \int_{-\pi}^{\pi} \cos(\cos(t)) \cos t\ dt - 0 \\
		& = 2i \int_{0}^{\pi} \cos(\cos(t)) \cos t\ dt \\
		& = 2i \int_{-\pi/2}^{\pi/2} \cos(\cos(u - \pi/2)) \cos(u - \pi/2)\ dt \\
		& = 2i \int_{-\pi/2}^{\pi/2} \cos(\sin(u)) \sin(u)\ dt \\
		& = 0
	\end{align*}
\end{solution}

\subquestion{}
\[
	\int_\gamma {\left(\frac{\Log(z)}{z}\right)}^2\ dz
\]
עבור $\gamma = e^{-it}$ ב־ $t \in [-\frac{\pi}{2}, \frac{\pi}{2}]$.
\begin{solution}
	הפעם $\gamma'(t) = -i e^{-it}$ ולכן
	\begin{align*}
		\int_\gamma {\left(\frac{\Log(z)}{z}\right)}^2\ dz
		& = \int_{-\frac{\pi}{2}}^{\frac{\pi}{2}} {\left(\frac{\Log(e^{-it})}{e^{it}}\right)}^2 \cdot (-i e^{-it})\ dt \\
		& = i \int_{-\frac{\pi}{2}}^{\frac{\pi}{2}} t^2 e^{-3it}\ dt \\
		& = {\left[ t^2 \frac{e^{-3it}}{-3}\right]}_{-\frac{\pi}{2}}^{\frac{\pi}{2}} + \frac{2}{3} \int_{-\frac{\pi}{2}}^{\frac{\pi}{2}} t e^{-3it}\ dt \\
		& = \frac{1}{2} \cdot \frac{1}{-3} + {\left[ t \frac{e^{-3it}}{-3i}\right]}_{-\frac{\pi}{2}}^{\frac{\pi}{2}} + \frac{2}{9i} \int_{-\frac{\pi}{2}}^{\frac{\pi}{2}} e^{-3it}\ dt \\
		& = \frac{1}{6} + 0 + \frac{1}{27 \cdot 2}
	\end{align*}
	זאת בשל חישוב ישיר, והעובדה ש־$e^{-\frac{\pi}{6}i} - e^{-\frac{\pi}{6}i} = \sin(-\frac{\pi}{6}) = -\frac{1}{2}$.
\end{solution}

\question{}
תהי $f : G \to \CC$ רציפה ותהי $\gamma : [a, b] \to G$ מסילה.

\subquestion{}
נסמן $M = \max_{t \in [a, b]} |f(\gamma(t))|$ ו־$L(\gamma)$ אורך המסילה $\gamma$. \\*
נראה כי מתקיים
\[
	\left\lvert \int_\gamma f(z)\ dz\right\rvert \le M \cdot L(\gamma)
\]
\begin{proof}
	\begin{align*}
		\left\lvert \int_\gamma f(z)\ dz \right\rvert
		& = \left\lvert \int_a^b f(\gamma(t)) \gamma'(t)\ dt \right\rvert \\
		& \le \int_a^b \left\lvert f(\gamma(t)) \gamma'(t) \right\rvert\ dt \tag{1} \\
		& = \int_a^b \left\lvert f(\gamma(t)) \right\rvert \cdot \left\lvert \gamma'(t) \right\rvert\ dt \\
		& \le \int_a^b M \cdot \left\lvert \gamma'(t) \right\rvert\ dt \tag{2} \\
		& = M \cdot \int_a^b \left\lvert \gamma'(t) \right\rvert\ dt \\
		& = M \cdot L(\gamma)
	\end{align*}
	כאשר
	\begin{enumerate}
		\item טענה ששאבנו מאינטגרלים ממשיים מרובי משתנים, ההוכחה עבור המקרה המרוכב זהה לחלוטין.
		\item מאינפי 3.
	\end{enumerate}
	ובכך הראינו כי אי־השוויון אכן חל.
\end{proof}

\subquestion{}
נוכיח כי לכל פונקציה על, מונוטונית עולה וגזירה ברציפות למקוטעין $\mu : [c, d] \to [a, b]$ מתקיים
\[
	\int_\gamma f(z)\ dz
	= \int_{\gamma \circ \mu} f(z)\ dz
\]
\begin{proof}
	נבחין תחילה ש־$\mu$ היא פונקציה הפיכה מההגדרה שלה. \\*
	נעבור לבחינת האינטגרל
	\[
		\int_{\gamma \circ \mu} f(z)\ dz
		= \int_c^d f(\gamma(\mu(t))) \cdot (\gamma \circ \mu)'(t)\ dt
		= \int_c^d f(\gamma(\mu(t))) \cdot \gamma'(\mu(t)) \cdot \mu'(t)\ dt
	\]
	אבל מכלל ההצבה האינטגרלי עבור $\mu^{-1}$ נובע
	\[
		\int_c^d f(\gamma(\mu(t))) \cdot \gamma'(\mu(t)) \cdot \mu'(t)\ dt
		= \int_a^b f(\gamma(\mu(\mu^{-1}(t)))) \cdot \gamma'(\mu(\mu^{-1}(t)))\ dt
		= \int_a^b f(\gamma(t)) \cdot \gamma'(t)\ dt
		= \int_\gamma f(z)\ dz
	\]
	כפי שרצינו.
\end{proof}

\subquestion{}
נניח ש־$f$ אנליטית ב־$G$ ונוכיח שמתקיים
\[
	\int_\gamma f'(z)\ dz
	= f(\gamma(b)) - f(\gamma(a))
\]
\begin{proof}
	על־ידי שימוש באינטגרל רימן, באדיטיביות האינטגרל ובתכונה שהוכחנו בסעיף א', נוכל לשחזר את ההוכחה של אינפי 2 לנוסחה היסודית הכללית. \\*
	בקצרה, נוכל לחלק את המסילה למספר הולך וגדל של תת־מסילות כך שנקודות הקצה הן היחידות שמשפיעות על גודל האינטגרל המלא, זאת תוך שימוש ברציפות הפונקציה ונגזרתה.
\end{proof}

\question{}
\subquestion{}
נמצא דוגמה לפונקציה אנליטית $f : G \to \CC$ ולשתי נקודות שונות $z_0, z_1 \in G$ כך שלכל $z \in G$ מתקיים
\[
	f'(z) \ne \frac{f(z_1) - f(z_0)}{z_1 - z_0}
\]
\begin{solution}
	נגדיר $f(z) = e^z$ ולכן $f'(z) = e^z$, ובנוסף נגדיר $z_1 = 2\pi i, z_0 = 0$ ולכן
	\[
		\frac{f(z_1) - f(z_0)}{z_1 - z_0} = \frac{1 - 1}{2\pi i - 0} = 0
	\]
	אבל אנו יודעים כי $e^z \ne 0$ לכל $z \in \CC$, ולכן נוכל להסיק כי $f'(z) \ne 0$ גם כן.
\end{solution}

\subquestion{}
תהי $f : G \to \CC$ פונקציה אנליטית ו־$z_0, z_1 \in G$ נקודות שונות כך ש־$[z_0, z_1] \subseteq G$. \\*
נראה שמתקיים
\[
	\left\lvert \frac{f(z_1) - f(z_0)}{z_1 - z_0} \right\rvert \le \max_{z \in [z_0, z_1]} |f'(z)|
\]
\begin{proof}
	נגדיר את המסילה $\gamma = [z_0, z_1]$, ולכן
	\[
		\left\lvert \int_\gamma f'(z)\ dz \right\rvert
		= |f(\gamma(1)) - f(\gamma(0))|
		= |f(z_1) - f(z_0)|
	\]
	מצד שני מ־ML נובע גם
	\[
		\left\lvert \int_\gamma f'(z)\ dz \right\rvert
		\le L(\gamma) \cdot \max_{z \in [z_0, z_1]} |f'(z)|
		= |z_1 - z_0| \cdot \max_{z \in [z_0, z_1]} |f'(z)|
	\]
	ולכן נובע משוויונות אלה
	\[
		|f(z_1) - f(z_0)|
		\le |z_1 - z_0| \cdot \max_{z \in [z_0, z_1]} |f'(z)|
		\iff
		\left\lvert \frac{f(z_1) - f(z_0)}{z_1 - z_0} \right\rvert \le \max_{z \in [z_0, z_1]} |f'(z)|
	\]
	כמבוקש.
\end{proof}

\question{}
יהיו $G, H \subseteq \CC$ תחומים ותהי $f : G \times H \to \CC$ פונקציה רציפה בשני משתנים. \\*
נאמר ש־$f(z, w)$ אנליטית ב־$z$ אם הנגזרת החלקית $\frac{\partial f}{\partial z}(z_0, w) = \lim_{z \to z_0} \frac{f(z, w) - f(z_0, w)}{z - z_0}$ קיימת ורציפה בכל נקודה בתחום. \\*
נגדיר עבור מסילה $\gamma : I \to H$ את הפונקציה $\varphi : G \to \CC$ על־ידי
\[
	\varphi(z) = \int_\gamma f(z, w)\ dw
\]

\subquestion{}
נקבע $z_0 \in G$ ונגדיר את הפונקציה $F_{z_0} : G \times H \to \CC$ על־ידי
\[
	F_{z_0}(z, w) = \begin{cases}
		\frac{f(z, w) - f(z_0, w)}{z - z_0} & z \ne z_0 \\
		\frac{\partial f}{\partial z}(z_0, w) & z = z_0
	\end{cases}
\]
נראה ש־$F_{z_0}$ רציפה כפונקציה בשני משתנים.
\begin{proof}
	בנקודות $z \ne z_0$ היא רציפה כהרכבת פונקציות רציפות, ולכן מספיק שנבדוק את המקרה $z = z_0$.
	\[
		\lim_{z \to z_0} F_{z_0}
		= \lim_{z \to z_0} \frac{f(z, w) - f(z_0, w)}{z - z_0}
		= \frac{\partial f}{\partial z}(z_0, w)  
		= F_{z_0}(z_0)
	\]
	ולכן היא אכן רציפה.
\end{proof}

\subquestion{}
נוכיח כי $\varphi$ אנליטית כך שמתקיים
\[
	\varphi'(z_0)
	= \int_\gamma F_{z_0}(z_0, w)\ dw
	= \int_\gamma \frac{\partial f}{\partial z}(z_0, w)\ dw
\]
\begin{proof}
	נבדוק את הגדרת הנגזרת
	\begin{align*}
		\lim_{z \to z_0} \frac{\varphi(z) - \varphi(z_0)}{z - z_0}
		& = \lim_{z \to z_0} \frac{\int_\gamma F_{z_0}(z, w)\ dw - \int_\gamma F_{z_0}(z_0, w)\ dw }{z - z_0} \\
		& = \lim_{z \to z_0} \int_\gamma \frac{F_{z_0}(z, w) - F_{z_0}(z_0, w)}{z - z_0}\ dw \\
		& = \lim_{z \to z_0} \int_\gamma F_{z_0}(z, w)\ dw \\
		& = \int_\gamma \lim_{z \to z_0} F_{z_0}(z, w)\ dw \\
		& = \int_\gamma F_{z_0}(z_0, w)\ dw \\
		& = \int_\gamma \frac{\partial f}{\partial z}(z_0, w)\ dw
	\end{align*}
	כאשר כלל המעברים נובעים משוויונות נתונים וחוקי אינטגרלים מרובי משתנים שאנו כבר מכירים. \\*
	בהתאם מרציפות האינטגרל והפונקציה $F_{z_0}$ נסיק כי $\varphi'$ אכן מוגדרת.
\end{proof}

\subquestion{}
נסיק כי $\psi(z) = \int_0^{2\pi} \cos(z e^{it})\ dt$ אנליטית.
\begin{proof}
	נגדיר $f(z, w) = \frac{\cos(zw)}{iw}$ וגם $\gamma : [0, 2\pi] \to \CC$ ךכ ש־$t \mapsto e^{it}$ ולכן
	\[
		\varphi(z)
		= \int_\gamma f(z, w)\ dw
		= \int_0^{2\pi} f(z, e^{it}) ie^{it}\ dt
		= \int_0^{2\pi} \frac{\cos(z e^{it})}{ie^{it}} ie^{it}\ dt
		= \psi(z)
	\]
	ולכן מסעיף ב' נובע ש־$\psi$ אכן אנליטית.
\end{proof}

\end{document}
